\documentclass[11pt]{article}
\usepackage[margin=1in]{geometry}
\usepackage{amsmath,amssymb,mathtools}
\usepackage{booktabs}
\usepackage{enumitem}
\usepackage{siunitx}
\sisetup{group-separator={,},group-minimum-digits=3}

\title{\textbf{Solutions -- Quiz 1} \vspace{ 0.4cm }\\ECON 300: Labour Economics}
\author{Instructor: Muhebullah Karimzada}
\date{January 22,  2026}

\begin{document}
\maketitle
\section*{Multiple Choice (5 points)}

\begin{enumerate}
\item \textbf{Answer: (c)}. A worker on temporary layoff expecting recall is counted as unemployed (they are without work but still attached to the labour market). A discouraged worker who stopped searching is \emph{not} counted as unemployed because they are not actively searching.

\item \textbf{Answer: (b)}. The LFPR is:
\[
LFPR = \frac{\text{Labour Force}}{\text{Working-age population}},
\]
and the labour force equals employed plus unemployed.

\item \textbf{Answer: (c)}. A higher wage raises the opportunity cost of leisure: one more hour of leisure means giving up more earnings.

\item \textbf{Answer: (b)}. The substitution effect makes leisure relatively more expensive, so individuals substitute away from leisure toward work, increasing hours worked.

\item \textbf{Answer: (c)}. A stronger preference for leisure increases the wage required to induce the person to give up leisure for work, raising the reservation wage.
\end{enumerate}


\subsection*{Problem 1: Labour Force Statistics and Interpretation (7.5 points)}

\textbf{Given:}
\[
\text{Working-age population} = 1{,}000{,}000,\quad
E = 640{,}000,\quad
U = 60{,}000,\quad
N = 300{,}000.
\]

\subsubsection*{(a) Labour force and LFPR (2.5 points)}

\textbf{Step 1: Compute the labour force.} By definition,
\[
LF = E + U = 640{,}000 + 60{,}000 = 700{,}000.
\]

\textbf{Step 2: Compute the labour force participation rate.}
\[
LFPR = \frac{LF}{\text{Working-age population}}
= \frac{700{,}000}{1{,}000{,}000} = 0.70 = 70\%.
\]

\subsubsection*{(b) Unemployment rate and employment-to-population ratio (2.5 points)}

\textbf{Step 1: Unemployment rate.}
\[
u = \frac{U}{LF} = \frac{60{,}000}{700{,}000}
= 0.085714\ldots \approx 8.57\%.
\]

\textbf{Step 2: Employment-to-population ratio.}
\[
\frac{E}{\text{Working-age population}} = \frac{640{,}000}{1{,}000{,}000}
= 0.64 = 64\%.
\]

\subsubsection*{(c) Recession with job losses and discouraged workers (2.5 points)}

\textbf{Shock 1: 20,000 employed lose jobs.} Those workers move from employed to unemployed (initially).
\[
E' = 640{,}000 - 20{,}000 = 620{,}000,
\qquad
U' = 60{,}000 + 20{,}000 = 80{,}000.
\]

\textbf{Shock 2: 15,000 discouraged workers stop searching.} This does not affect the labour force.  They were not part of labour force and still they are not part of it.

\[
U'' = 80{,}000 
\]

Since they exit the labour force, the new labour force becomes:
\[
LF'' = E' + U'' = 620{,}000 + 80{,}000  = 700{,}000.
\]

\textbf{New unemployment rate:}
\[
u'' = \frac{U''}{LF''} = \frac{80{,}000 }{700{,}000}
= 0.11429\ldots \approx 11.43\%.
\]


\newpage
\subsection*{Problem 2: Labour Supply, Market Equilibrium, and Minimum Wage (7.5 points)}

\textbf{Given individual problem:}
\[
C = 600 + 25h,
\qquad
U = C - 0.5h^2.
\]

\textbf{Given market:}
\[
L^D = 120 - 2w,
\qquad
L^S = 20 + 3w.
\]

\subsubsection*{(a) Optimal hours worked (2.5 points)}

\textbf{Step 1: Substitute the budget constraint into utility.}
\[
U(h) = (600 + 25h) - 0.5h^2
= 600 + 25h - 0.5h^2.
\]

\textbf{Step 2: Differentiate with respect to $h$ and set equal to zero (FOC).}
\[
\frac{dU}{dh} = 25 - h.
\]
Set FOC:
\[
25 - h = 0 \quad \Rightarrow \quad h^* = 25.
\]

\textbf{Step 3: Check second derivative (SOC).}
\[
\frac{d^2U}{dh^2} = -1 < 0,
\]
so $h^*=25$ is a maximum.

\textbf{Answer:} The worker supplies \textbf{25 hours per week}.

\subsubsection*{(b) Market equilibrium wage and employment (2.5 points)}

\textbf{Step 1: Set supply equal to demand.}
\[
120 - 2w = 20 + 3w.
\]

\textbf{Step 2: Solve for $w$.}
\[
120 - 20 = 3w + 2w
\quad \Rightarrow \quad
100 = 5w
\quad \Rightarrow \quad
w^* = 20.
\]

\textbf{Step 3: Substitute back to get equilibrium employment.}
Using demand:
\[
L^* = 120 - 2(20) = 120 - 40 = 80.
\]
(Check with supply: $20 + 3(20)=20+60=80$.)

\textbf{Answer:} Equilibrium wage is \textbf{$w^*=20$} and employment is \textbf{$L^*=80$}.

\subsubsection*{(c) Minimum wage $w=20$ (2.5 points)}

\textbf{Step 1: Compare minimum wage to equilibrium wage.}
\[
w_{\min} = 20,\quad w^* = 20.
\]
Because the minimum wage equals the competitive equilibrium wage, it is \textbf{not binding}.

\textbf{Step 2: Compute labour demanded and supplied at $w=20$.}
\[
L^D(20) = 120 - 2(20) = 80,
\qquad
L^S(20) = 20 + 3(20) = 80.
\]

\textbf{Step 3: Unemployment (excess supply).}
\[
U = L^S - L^D = 80 - 80 = 0.
\]

\textbf{Answer:} The minimum wage is \textbf{not binding} and creates \textbf{zero unemployment}.

\end{document}
